\section{System implementation}
\label{sec:system}

\subsection{Coordination protocols}

\begin{figure}[H]
\centering
\includegraphics[width=0.55\linewidth,height=0.80\textheight,keepaspectratio]{topologies.png}
\caption{The four coordination topologies.}
\label{fig:schematic-topologies}
\end{figure}

Table~\ref{tab:topology-protocol} summarises the four protocols at the sweep's fixed
setting of three agents and two rounds.

\intable{tab_topology_protocol}

\textbf{Single agent.} One agent answers once with an empty peer context. The result
defines the baseline used by every efficiency metric: one turn, zero messages, one round.

\textbf{Independent.} Three agents answer in parallel with empty peer context and
distinct seeds, and the system answer is the majority vote over their outputs. Ties are
broken by the summed option-logprob probability mass assigned to each tied answer, and if
that is also tied, by position. There is no communication, so this condition isolates the
ensemble effect of aggregation from the effect of information exchange.

\textbf{Decentralized.} Three agents debate all-to-all for two rounds. In round zero the
peer context is empty. In round one each agent receives the rendered outputs of both
other agents from round zero. The system answer is the majority vote over the final
round, not over all rounds. Message accounting begins at round one, so a three-agent
two-round cell records six inter-agent messages, which is $n(n-1)$ for one communicating
round.

\textbf{Centralized.} One agent acts as orchestrator and the remaining two act as
sub-agents. In each round the sub-agents answer in parallel and the orchestrator then
produces a synthesis conditioned on all sub-agent outputs at the cell's context-sharing
level. From round one onward the sub-agents see only the orchestrator's previous
synthesis, not each other. The system answer is the final synthesis, not a vote.

Two features of the centralized protocol deserve emphasis because they are easy to
misread. A three-agent centralized cell has only two answering sub-agents, so its vote
statistics are computed over two outputs rather than three. And its final answer is
produced by a single designated agent, which means the final-producer confidence signal
for centralized cells is the orchestrator's own probability, whereas for the two voting
topologies it is the probability of the decisive agent on the winning side.

Agent seeding is deterministic and structured so that no two agents in a cell share a
sampling stream. Independent agents use the cell seed offset by agent index. Debating
agents use the cell seed offset by round times one hundred plus agent index. The
orchestrator uses a reserved offset. Every generation call therefore has a reproducible
seed derived from the cell identifier.

\subsection{Context sharing}

Peer context is rendered by a single component with a frozen protocol hash, so that a
change to the rendering rules is detectable in the data rather than silent.

At the artifact-only level a peer block contains the peer's identifier, its chosen answer
and, when present, its stated confidence percentage. At the plus-CoT level the block
additionally contains the peer's chain of thought, truncated to 4{,}000 characters. An
intermediate level exists in the code, adding post-answer scratchpad content, but was
removed from the grid by the trim and does not occur in the analysis population.

The chain of thought is the model's native thinking trace when the serving stack returns
one, and otherwise the portion of the response preceding the final-answer marker. The
whole rendered peer block is additionally capped at 4{,}000 exact tokens under the serving
profile's tokenizer. When truncation occurs, the block carries an explicit marker
recording the original token count and a digest of the untruncated text, and each
consuming generation records the peer-context token count, its digest, the per-block token
counts and the number of truncation markers it saw.

The levels are nested by eligibility rather than by literal text inclusion: because the
token cap binds after the eligibility rules are applied, a plus-CoT block is not
guaranteed to be a textual superset of the artifact-only block it would otherwise have
been. A coordinated topology fails closed if the exact serving-profile tokenizer is
unavailable, so a cell never silently falls back to an approximate cap.

\subsection{Prompt templates}

The two executed prompt levels are reproduced verbatim in Appendix~\ref{app:axes}. Level 0
is a single sentence. Level 3 is a six-stage protocol whose stages are comprehension,
retrieval, reasoning, alternative consideration, checking and calibration, followed by an
explicit output-format block and a worked format example. The calibration stage instructs
the model to judge its confidence from the strength of its reasoning rather than defaulting
to high certainty.

Because prompt length and prompt quality are confounded by construction, both are measured
and stored per cell. Length is the exact token count under the model's tokenizer. Quality
is a deterministic composite of five binary presence features (role, strategy, format,
examples, constraints), the count of explicit numbered steps, a readability score and word
count, combined with fixed weights. An optional model-judged rubric score is also
supported. Recording both means a prompt effect can be attributed to structure rather than
merely to a longer context, though the two-level design cannot fully separate them.

Production workers read prompts from an immutable release worktree rather than from the
source checkout, so an edit to a template cannot retroactively change what a running cell
sends.

\subsection{Reasoning budget mechanism}

\begin{figure}[H]
\centering
\includegraphics[width=0.62\linewidth,height=0.52\textheight,keepaspectratio]{reasoning_ece_protocol.png}
\caption{Reasoning and confidence-measurement protocol. The deterministic option probe
runs before any stochastic generation, so reasoning depth cannot influence the
confidence measurement path.}
\label{fig:schematic-reasoning}
\end{figure}

Thinking budgets are enforced by the serving engine at the token level within a single
chat request \citep{kwon2023vllm}. They are not implemented by decoding, truncating and
re-encoding a continuation, which would change the model's own context. The protocol is
frozen under a version number and a specification hash covering the endpoint, the budget
field name, the reasoning parser, the required engine version, the sampling parameters,
the number of generation phases and the token allowances.

Table~\ref{tab:reasoning-protocol} states the rung mapping and the resulting output
floors.

\intable{tab_reasoning_protocol}

A request must be able to obtain its floor: 4{,}096 tokens reserved for the answer, plus
the thinking allowance for the rung. The actual token limit sent is the whole remaining
envelope, that is the served context length minus the exact local prompt token count minus
a 128-token reserve, and the request is required to reach the served context boundary
exactly. A cell whose rendered prompt cannot leave its floor within the standard serving
profile is routed deterministically to a long profile with a 40{,}960-token context, and
the two profiles expose the same served model name so that runtime layout never leaks into
the capacity treatment.

Compliance is verified after the fact rather than assumed. Each completion is checked
against the delimiter contract: a thinking completion must open its trace at position zero
and close it later; the closing position must not exceed the requested budget; and the
terminal token must be the expected end-of-message token. Violations are recorded as
typed codes from a frozen set, and a violating response is retained as censored data
rather than resampled, so that protocol failures cannot bias the sample by selective
regeneration. Cells containing any censored question have their calibration statistics set
to missing rather than computed on a filtered subset. The engine's reported prompt token
identities are also required to equal the locally rendered chat template exactly.

\subsection{Models and serving}

\begin{figure}[H]
\centering
\includegraphics[width=\linewidth,height=0.34\textheight,keepaspectratio]{execution_architecture.png}
\caption{Execution architecture: scheduler fleet, serving endpoints and the resumable
cell dispatcher.}
\label{fig:schematic-exec}
\end{figure}

Table~\ref{tab:models} lists the served models.

\intable{tab_models}

All six are dense Qwen3 checkpoints \citep{qwen3} pinned to explicit forty-character
commit revisions, with tokenizer identifier and revision pinned identically and a
checksum sidecar over the contract file. Weights are loaded in bfloat16 at 0.90 GPU memory
utilisation with a twenty-logprob cap and the reasoning parser enabled, on A100-80GB
hardware, with the engine version asserted inside each serving job. Dataset, model and
tokenizer downloads are disabled at run time, so a job either finds its pinned artefacts
in the local cache or fails.

The 32B model is served at 16{,}384 tokens of context rather than 32{,}768. This is a
memory constraint: 32B weights in bfloat16 occupy roughly 61 GiB, which leaves
insufficient key-value cache room for the full context on a single A100. The consequence
is a genuine confound in one corner of the design, and it is flagged per cell and analysed
explicitly in Section~\ref{sec:health-cap}.

Table~\ref{tab:sampling} states the sampling configuration.

\intable{tab_sampling}

Two details matter for interpreting the calibration results. First, the sampling
temperature differs between the thinking-enabled and thinking-disabled conditions, at 0.6
and 0.7 respectively, with correspondingly different nucleus thresholds. The reasoning
axis is therefore not a pure single-variable manipulation. Second, the confidence probe is
a separate deterministic request at temperature zero with a single generated token, issued
before any stochastic generation for that agent and cached, so that a later probe failure
cannot cause a chat response that has already been observed to be reissued.

\subsection{Benchmarks, question contracts and grading}

\begin{figure}[H]
\centering
\includegraphics[width=\linewidth,height=0.30\textheight,keepaspectratio]{per_question_flow.png}
\caption{Per-question execution flow, from question rendering through the deterministic
option probe and the topology run to grading and record assembly.}
\label{fig:schematic-flow}
\end{figure}

Table~\ref{tab:benchmarks} lists the evaluation suites.

\intable{tab_benchmarks}

Each suite is pinned to a forty-character dataset commit revision rather than a mutable
branch, and is loaded cache-only with the revision verified against the on-disk path of
the Arrow file that was actually read. The per-cell question count is the minimum of the
requested 200 and the split size, which is why GPQA Diamond contributes 198 questions and
the other three contribute 200. Two narrow direct-Arrow fallbacks exist for metadata
incompatibilities in two of the suites; both run the same normalisation and record a
checksum of the file they read.

All suites are normalised to a common question record. Multiple-choice options are
shuffled deterministically per question using a seed derived from the cell seed and the
question index, so the correct answer is not systematically in the first position and the
shuffle is reproducible. Option letters are assigned by position after shuffling.

Before a run begins, the exact question list for each combination of benchmark, question
count and seed is frozen into a contract file recording an ordered list of question
identifiers, a per-question digest, a digest over the whole question contract and a digest
of the source identity. The contract is bound to the cell manifest by its own hash. At run
time, questions are served through a verified catalogue that fails closed on any mismatch,
so a silent dataset change cannot alter what a resumed cell evaluates.

Questions are rendered to a fixed format. Multiple-choice questions present the stem, the
lettered options and an instruction to respond with a letter. Free-form questions instruct
the model to place its final answer in a boxed expression. Every question, in both
families, appends an instruction to state a confidence percentage on a new line after the
answer. When thinking is disabled, a short chain-of-thought hint is prepended to the user
message; when thinking is enabled it is not, because the native thinking mode supplies the
same function.

Grading extracts a letter using four ordered regular expressions with a bare-letter
fallback, or extracts the last brace-balanced boxed expression and normalises it by
stripping formatting commands, whitespace, currency and percent symbols and trailing
punctuation. If chat parsing fails on a multiple-choice question, the answer falls back to
the argument maximum of the option probe, so a formatting failure does not silently become
an incorrect answer.

\subsection{Execution}

Cells are dispatched as independent array tasks against a fleet of serving endpoints under
a shared scheduler. The dispatcher is resumable at question granularity: a cell that is
preempted resumes by question identifier and re-runs only what is missing. Serving
endpoints are maintained by a keepalive process against a declared fleet specification and
are re-registered rather than duplicated when a job survives a controller restart.

This execution model has one direct consequence for the data. Because cells append to their
result streams and can be resumed after partial writes, duplicate question rows are
possible, and malformed trailing lines are possible when a task is killed mid-write. Both
occur in this run and both are handled at ingestion time, quantified in
Section~\ref{sec:health}.
